\section{An explicit circuit joining the three constructions}\label{sec:circuit}
The logarithmic coefficient in Section~\ref{sec:resonance} predicts a shear of the solution basis. We can produce the continuation and its nonzero increment for an actual loop. This links a finite path presentation, quantitative avoidance of a singular point, and the analytic coefficient calculation in one example.

\begin{theorem}[Monodromy of the resonant basis on a polygon]\label{thm:circuit}
In the notation of Theorem~\ref{thm:res-gauss}, put $\rho=R/4$. Use the positively oriented diamond in the $t=1-z$ plane through $\rho,i\rho,-\rho,-i\rho,\rho$, with the real logarithm at the starting point. One turn fixes $U$ and takes $V$ to $V+\Lambda g_kU$, where
\[
 \Lambda=4i\int_0^1\frac{ds}{(1-s)^2+s^2}
         =8i\arctan(1)=2\pi i .
\]
There is a PR name for this transport from the rational parameter data alone. For $\nu\in\mathbb Z$, the $\nu$-fold circuit adds $\nu\Lambda g_k U$ to $V$. The cyclic action is trivial exactly when $(a+s)_k(b+s)_k=0$, with $s=\min(0,m)$ as in the local theorem. Otherwise every nonzero power acts nontrivially.
\end{theorem}

\begin{proof}
Along the $j$th side write $t(u)=\rho i^j(1-u+iu)$, $0\le u\le1$. Its $\ell^1$ norm is exactly $\rho$ and $|1/t(u)|_1\le2/\rho$. Divide each side into sixteen equal subsegments. From any initial vertex $p$ of a subsegment the displacement $h'$ within it satisfies $|h'/p|_1\le1/4$. On this subsegment transport the logarithm by the convergent increment
\[
 \sum_{n=1}^{N}\frac{(-1)^{n-1}}{n}(h'/p)^n,\qquad
 \text{tail }\ \le \tfrac43\,4^{-N-1}.
\]
This defines a branch near every subsegment, agrees with the transported value at its starting vertex, and has derivative $1/t$. Indeed the derivative of the finite sum is the partial geometric series for $1/(p+h')$, whose remainder is bounded uniformly on the segment. This gives all the branch data used in the local construction. In the other terms of $H$, integer powers and the series for $g$ return to their initial values. Uniform convergence follows from the weights on $|t|_1=\rho$.

For a direct computation of the logarithm increment on an edge, its derivative with respect to $u$ is
\[
 \frac{-1+i}{1-u+iu}=\frac{2u-1+i}{D(u)},\qquad
 D(u)=(1-u)^2+u^2 .
\]
The increments just constructed sum to the integral of this rational function. At finite precision this follows by integrating the finite geometric polynomial on each subsegment and bounding the integral of its remainder by the subsegment length times the supremum error. Every such integral of a polynomial uses rational coefficients and integer divisors. Also $D\ge 1/2$, $|D'|\le2$. The real part cancels exactly in midpoint sums when $u$ is paired with $1-u$. The imaginary part $D^{-1}$ lies between $1$ and $2$ and is $8$-Lipschitz, since
$|D(u)-D(v)|\le2|u-v|$ and $D(u)D(v)\ge1/4$.
A uniform midpoint rule with $M$ intervals has error at most $4/M$ on one side by comparison on each interval; four sides cost at most $16/M$. These comparisons are statements about finite sums, since any finer sum differs from the chosen midpoint integral by that bound, or by the same bound plus the finer mesh error using a common refinement. This constructs the displayed increment with $\Im\Lambda\ge4$. Substitution $v=2u-1$ and the evenness of $1/(1+v^2)$ give $2\int_0^1(1+v^2)^{-1}dv=2\arctan(1)$ per side in the imaginary part. One can use $\pi=4\arctan(1)$ as its name convention. The lower bound already proves apartness without any identity about $\pi$.

The change of $H$ is $\Lambda g_k$, hence the change of $V=UH$ is the stated shear. Reversing a segment reverses its integral. Induction on a finite number of positive or negative traversals yields the factor $\nu$. To read the induction in iPRA, store the signed traversal count and the multiplier of the common named constant $\Lambda g_k$; the update is an exact signed-integer addition. To evaluate the transported jet, approximate the constant once with the additional precision required for multiplication by $|\nu|$. When $g_k\ne0$, its inverse is explicit Gaussian rational, $\Lambda$ is apart from zero, and $U(\rho)=\rho^r h(\rho)$ in the notation of Section~\ref{sec:resonance} has a named inverse. Thus a nonzero power changes the basis at the initial point itself. The rising factorial test follows from the closed formula for $g_k$ and decidable rational arithmetic.
\end{proof}

At $(a,b,c)=(1,1,3)$ the logarithmic coefficient is $-1$. The regular solution from zero has local expression $-2U-2V$, so one positive circuit changes it by $2\Lambda U$. The finite-presentation statement here concerns the cyclic group generated by the displayed polygon. A claim about every path with a prescribed winding number would require a route from such a path to this presentation, for example a validated homotopy or a direct integral calculation.
